\documentclass[11pt,letterpaper]{article}
\usepackage[T1]{fontenc}
\usepackage[utf8]{inputenc}
\usepackage{lmodern,microtype}
\usepackage[margin=1in]{geometry}
\usepackage{amsmath,amssymb,amsthm,mathtools}
\usepackage{booktabs,longtable,array,enumitem}
\usepackage{xcolor,listings,fancyhdr}
\usepackage[hyphens]{url}
\usepackage[colorlinks=true,linkcolor=blue!45!black,citecolor=blue!45!black,urlcolor=blue!45!black]{hyperref}
\usepackage[nameinlink,noabbrev]{cleveref}
\setlength{\emergencystretch}{3em}
\setlength{\parskip}{3pt}
\setcounter{tocdepth}{1}
\setlist{itemsep=3pt,topsep=5pt}
\pagestyle{fancy}
\fancyhf{}
\fancyhead[L]{\small Pell equations in iPRA}
\fancyhead[R]{\small Explicit finite constructions}
\fancyfoot[C]{\thepage}
\setlength{\headheight}{14pt}
\newtheorem{theorem}{Theorem}[section]
\newtheorem{lemma}[theorem]{Lemma}
\newtheorem{proposition}[theorem]{Proposition}
\newtheorem{corollary}[theorem]{Corollary}
\theoremstyle{definition}
\newtheorem{definition}[theorem]{Definition}
\newtheorem{remark}[theorem]{Remark}
\newcommand{\N}{\mathbb N}
\newcommand{\Z}{\mathbb Z}
\newcommand{\Q}{\mathbb Q}
\newcommand{\ipra}{\mathsf{iPRA}}
\newcommand{\isqrt}{\operatorname{isqrt}}
\newcommand{\Nm}{\operatorname{N}}
\newcommand{\CFB}{\mathcal B}
\newcommand{\collB}{\mathcal H}
\newcommand{\pair}[2]{(#1,#2)}
\newcommand{\abs}[1]{\lvert#1\rvert}
\newcommand{\code}[1]{\texttt{#1}}
\lstset{basicstyle=\small\ttfamily,columns=fullflexible,breaklines=true,
  frame=single,rulecolor=\color{black!20},backgroundcolor=\color{black!2},
  showstringspaces=false,keepspaces=true}
\hypersetup{pdftitle={Pell Equations in Intuitionistic Primitive Recursive Arithmetic},
  pdfauthor={OpenAI},pdfsubject={Constructive number theory with explicit primitive recursive bounds}}

\begin{document}
\begin{titlepage}
\vspace*{1.1cm}
{\Large\color{blue!45!black} Constructive number theory\par}
\vspace{0.5cm}
{\huge\bfseries\raggedright Pell Equations in\\
Intuitionistic Primitive\\
Recursive Arithmetic\par}
\vspace{0.7cm}
{\Large\raggedright Explicit finite constructions, complete solution theory,\\
and a Lean formalization roadmap\par}
\vspace{1cm}
{\large Prepared for Vladimir Reshetnikov\par}
{\large November 15, 2026\par}
\vfill
\begin{abstract}
We develop the ordinary and generalized Pell equations using only intuitionistic
logic, primitive recursive functions, and quantifier-free induction. For every
nonsquare integer $D\ge2$, we construct its fundamental positive solution,
classify all solutions, decide the negative equation and every equation
$x^2-Dy^2=N$, and prove divisibility, congruence, and approximation theorems.
The existence argument is given twice: first by two explicitly finite
pigeonhole arguments, then by a reduced continued-fraction algorithm with at
most $4a(a+1)$ convergents, where $a=\isqrt(D)$. A finite proof of Legendre's
criterion establishes the familiar period-parity rule and minimality.
All searches, descents, and approximation moduli have stated bounds.
Radical notation denotes a decidable ordered field of rational pairs, not an
assumed real-number continuum. The accompanying code exercises the exact
arithmetic constructions. A detailed plan connects the proofs to the actual
embedded iPRA theory and native-calculus bridge in ProveIt.
\end{abstract}
\vspace{0.4cm}
\noindent\textbf{Formalization status.} This is a fully developed human-readable
mathematical argument with an explicit iPRA induction audit. It is not a
completed machine-checked iPRA proof. The supplied Lean algebraic starting
point and transfer example were not compiled in this environment. The
remaining object-theory proof obligations are listed explicitly.
\end{titlepage}

\tableofcontents
\newpage

\section{Scope and principal results}\label{sec:scope}

Throughout the main development, $D$ is a natural number satisfying
\begin{equation}\label{eq:guard}
 D\ge2,\qquad \isqrt(D)^2\ne D.
\end{equation}
The integer $D$ need not be squarefree. This matters for both the negative
equation and the choice of quadratic order: the coefficients of our solutions
are integers in $\Z[\sqrt D]$, and we do not replace this order by the ring of
integers of a quadratic number field.

The following is a precise finite form of the theory developed below.

\begin{theorem}[Constructive Pell theory]\label{thm:main}
There are fixed primitive recursive algorithms with the following properties,
whose correctness has proofs using intuitionistic logic and quantifier-free
induction. Put
\[
 a=\isqrt(D),\qquad S=2a(a+1),\qquad
 \CFB(D)=(2a+1)^{2S}.
\]
\begin{enumerate}
\item The algorithm for the fundamental solution returns positive integers
$(A,B)$ with $A^2-DB^2=1$, $A,B\le\CFB(D)$, and $B$ least among positive
second coordinates of solutions. It examines at most $2S=4a(a+1)$ continued-
fraction convergents.
\item Define $(X_0,Y_0)=(1,0)$ and
\[
 (X_{n+1},Y_{n+1})=(AX_n+DBY_n,\ BX_n+AY_n).
\]
Every nonnegative norm-one solution is exactly one $(X_n,Y_n)$. Its exponent is
computed by a descent with at most its input second coordinate many active
steps. Independent changes of coordinate signs give every integer solution.
\item The continued-fraction period $\ell$ is at most $S$. Negative Pell is
solvable exactly when $\ell$ is odd. In that case a computed positive solution
$v$ of norm $-1$ satisfies $v^2=(A,B)$, and all positive negative-Pell
solutions are $v(A,B)^n=v^{2n+1}$.
\item For any nonzero integer $N$, let $K=\abs N$. Solvability of
$x^2-Dy^2=N$ is equivalent to the existence of a solution in the finite box
\[
 0\le y\le AK,\qquad 0\le x\le(DA+1)K.
\]
Every nonnegative solution is a member of this box multiplied by a power of
$(A,B)$. Thus generalized Pell solvability is uniformly primitive recursive.
For $N=0$, the only solution is $(0,0)$.
\item The power sequence satisfies
\[
 \gcd(Y_m,Y_n)=Y_{\gcd(m,n)}.
\]
For each $h\ge1$ a return to $(1,0)$ modulo $h$ is computable within $h^2$
steps, and the residue sequence is periodic from its first term.
\item The rational intervals
\[
 \left[\frac{DY_n}{X_n},\frac{X_n}{Y_n}\right]\quad(n\ge1)
\]
are nested, straddle the positive square root in the sense of squared
inequalities, and have width at most $2^{-2n+2}$. The upper endpoints are
one-sided best approximants among rational upper approximations with no
larger denominator.
\end{enumerate}
\end{theorem}

``Primitive recursive'' describes a uniform arithmetic construction, not a
claim of practical speed in the binary length of $D$. The bound $\CFB(D)$ is
deliberately loose; its role is to make totality explicit. The continued-
fraction implementation is far more useful computationally than scanning the
whole coordinate box.

The standard number-theoretic background is discussed in Conrad's two
expositions \cite{conrad1,conrad2}. Our purpose is to give all the finite
proofs and the arithmetic bounds needed for the requested foundation. No
result from those expositions is used as an unproved mathematical premise.

\section{The arithmetic setting}\label{sec:arithmetic}

\subsection{Exactly which theory is being used}

We use first-order intuitionistic primitive recursive arithmetic: natural
numbers, function symbols for primitive recursive functions, their defining
equations, successor separation and injectivity, and induction restricted to
quantifier-free formulas. Ordinary first-order intuitionistic inference rules
are available. In particular, we do not assume induction for arbitrary
formulas containing unbounded quantifiers.

Integer and rational notation is an abbreviation for numerical codes.
For example, $2n$ codes $n\ge0$ and $2n-1$ codes $-n$ for $n\ge1$.
Rationals are signed numerators with positive denominators; cross
multiplication suffices for equality and order. Finite sequences can be
encoded by pairing or by positional notation with a sufficiently large
explicit base. Their access, update, and finite folds are primitive recursive.
No arbitrary sequence is quantified over as a second-order object.

For a fixed primitive recursive Boolean predicate $P(i,\vec z)$, the bounded
conjunction
\[
 \bigwedge_{i<t}P(i,\vec z)
\]
is represented by an iterated Boolean product. Its assertion is an equation
between primitive recursive terms. Bounded search is implemented by a
recursion that stores the first successful index and a success flag; it is
total even when no index succeeds. These devices explain the bounded
quantifiers used in the prose. The compilation and induction audit in
\cref{sec:audit} specifies their use in each substantial proof.

\subsection{Square roots, gcd, and rational irrationality}

Define $\isqrt(t)$ by testing the integers $0,\ldots,t$ and retaining the
largest whose square is at most $t$. The defining property is
\begin{equation}\label{eq:isqrt}
 \isqrt(t)^2\le t<(\isqrt(t)+1)^2.
\end{equation}
The upper inequality follows because the next trial would otherwise have
been accepted; the case $t=0$ is immediate. This is a bounded-search proof.
An optimized integer-square-root implementation is an implementation choice,
not an additional mathematical principle.

We will also use the Euclidean algorithm with a fixed supply of steps.
For nonnegative input $(r,s)$, a nonterminal remainder step strictly decreases
the second coordinate, since $r\bmod s<s$ when $s>0$.
Consequently $r+s+1$ stages suffice if the
terminal state is held fixed. The invariants
\[
 \gcd(r,s)=\gcd(s,r\bmod s)
\]
and the usual two accumulated linear combinations prove gcd and Bezout's
identity by quantifier-free induction. For clarity, the coprimality
consequence we need can be proved directly: if $ur+vs=1$, squaring the
identity gives
\[
 u^2r^2+s(2uvr+v^2s)=1.
\]
Thus a positive $s$ dividing $r^2$ must be $1$ whenever $r,s$ are coprime.

\begin{lemma}[No rational square root]\label{lem:irrational}
Under \eqref{eq:guard}, $p^2-Dq^2\ne0$ for integers $p$ and $q\ne0$.
\end{lemma}
\begin{proof}
Reduce $\abs p/\abs q$ by their gcd to $r/s$ with $s>0$ and
$\gcd(r,s)=1$. If $r^2=Ds^2$, then $s\mid r^2$, so the preceding Bezout
calculation gives $s=1$. It follows that $D=r^2$, contradicting
\eqref{eq:guard}. Every reduction and divisibility check here has the
fixed-fuel Euclidean implementation just described.
\end{proof}

\subsection{Radical notation without real numbers}\label{sec:quadraticorder}

It is convenient to write $s=\sqrt D$. This notation will refer to the
formal pair $(0,1)$ in a field of rational pairs, not to a real-number object.
Define
\[
 (u,v)+(r,t)=(u+r,v+t),\qquad
 (u,v)(r,t)=(ur+Dvt,ut+vr).
\]
Identify $(u,v)$ with $u+vs$, so $s^2=D$. Conjugation sends $u+vs$ to
$u-vs$, and the norm is $u^2-Dv^2$. By \cref{lem:irrational}, a nonzero
pair has nonzero norm after clearing rational denominators. Its inverse is
\[
 (u+vs)^{-1}=\frac{u-vs}{u^2-Dv^2}.
\]
The field identities follow by expansion of rational polynomials.

Here is a fully algebraic construction of its order. Let $\mathcal A$
consist of pairs with nonnegative rational coefficients, not both zero.
The set $\mathcal A$ is closed under addition and multiplication, and does
not contain zero. Declare $z>0$ when $z=P/Q$ for some $P,Q\in\mathcal A$.
This positive cone is closed under sums and products, using
\[
 \frac P Q+\frac R T=\frac{PT+RQ}{QT},\qquad
 \frac P Q\frac R T=\frac{PR}{QT}.
\]
The number $-1$ cannot lie in it: $P/Q=-1$ would imply $P+Q=0$, impossible
for $P,Q\in\mathcal A$.
The cone is also closed under inversion, by interchanging $P$ and $Q$.

The definition is decidable, with explicit witnesses rather than an
unbounded search for $P,Q$. If the two coefficients of a nonzero $u+vs$
have the same sign, the sign is immediate. For opposite signs, write for
example $u-bs$ with $u,b>0$ and use
\begin{equation}\label{eq:signformula}
 u-bs=\frac{u^2-Db^2}{u+bs}.
\end{equation}
The denominator belongs to $\mathcal A$, and the numerator is a nonzero
rational. This explicitly places exactly one of $u-bs$ and its negative
in the positive cone. The case $-u+bs$ is the negative of this case.
Zero coefficients are handled directly. If both $z$ and $-z$ were positive,
their quotient would make $-1$ positive, which was excluded. Thus this gives
a decidable total ordered field with $s>0$.

All the witnesses and sign cases above are rational arithmetic formulas.
Consequently this order construction can be compiled into iPRA using
finite sign distinctions and the norm calculation. It does not assert the
existence of an arbitrary ordered field as a set inside iPRA. In particular,
for $r\ge0$,
\[
 r<s\ \Longleftrightarrow\ r^2<D,
\]
and \eqref{eq:isqrt} gives $a<s<a+1$ when $a=\isqrt(D)$.
Every subsequent calculation involving $s$ takes place in this coded field.

\subsection{A witness-producing finite pigeonhole lemma}\label{sec:pigeonhole}

\begin{lemma}\label{lem:pigeonhole}
Given a coded list $v_0,\ldots,v_R$ whose entries are less than $R$, with
$R\ge1$, a bounded scan returns indices $i<j\le R$ with $v_i=v_j$.
Its correctness is provable with quantifier-free induction.
\end{lemma}
\begin{proof}
Keep an array of $R$ occupancy bits, initially zero, together with the first
index stored at each occupied position. Scan the list. On an unoccupied
position set its bit to one; on an occupied position return the stored index
and the current index. A success flag freezes an already returned answer.
After $t$ unsuccessful entries the sum of the occupancy bits is $t$:
each step changes exactly one zero bit to one. This sum is at most $R$
because each of its $R$ summands is at most one. Therefore $R+1$
unsuccessful entries would imply $R+1\le R$.

The last argument is constructive. The scan's success flag is decidable.
Its zero case gives the displayed arithmetic contradiction; its nonzero
case supplies the stored witnesses. The bit sum, the statement that all
bits are zero or one, and the stored-index invariants are primitive
recursive finite folds. Induction is applied to equations asserting these
invariants, not to a formula with an unbounded quantifier.
\end{proof}

\section{Existence from two finite collisions}\label{sec:collision}

This first proof makes the elimination of an infinite pigeonhole argument
particularly transparent. The subsequent continued-fraction proof gives a
better numerical procedure and a smaller bound.

\subsection{Exact finite rational approximation}

\begin{lemma}[Finite Dirichlet construction]\label{lem:dirichlet}
For $Q\ge1$ a bounded algorithm produces positive integers $p,q$ with
\begin{equation}\label{eq:dirichlet}
 1\le q\le Q,\qquad 1\le p\le Dq,\qquad
 \abs{p-qs}<\frac1Q.
\end{equation}
\end{lemma}
\begin{proof}
For $0\le j\le Q$ compute
\[
 t_j=\isqrt(Dj^2),\quad
 w_j=\isqrt(DQ^2j^2),\quad
 b_j=w_j-Qt_j.
\]
These are exact integer computations. The square-root inequalities give
$Qt_j\le w_j<Q(t_j+1)$, so $0\le b_j<Q$.
By \cref{lem:pigeonhole}, choose $i<j$ with $b_i=b_j$ and set
$q=j-i$, $p=t_j-t_i$.

To check the error, put $e_j=js-t_j$ in the coded ordered field. We have
$0\le e_j<1$ and $b_j\le Qe_j<b_j+1$. Equal bins therefore give
\[
 \abs{qs-p}=\abs{e_j-e_i}<1/Q.
\]
Because $q\ge1$ and $s>1$, $p>qs-1>0$. Also $s<D$, so
$p<qs+1\le Dq+1$, whence the integer $p$ is at most $Dq$.
This proves all of \eqref{eq:dirichlet}. The apparent fractional parts in
the verification are abbreviations for the two displayed integer square
roots; the algorithm never operates on an arbitrary real number.
\end{proof}

\subsection{Separated approximants and signatures}

Set
\begin{equation}\label{eq:collisionconstants}
 C=2D+1,\qquad M=2C^3,\qquad Q_i=C^{i+1}\quad(0\le i\le M).
\end{equation}
For each of these $M+1$ explicitly specified inputs run
\cref{lem:dirichlet}, obtaining $p_i,q_i$, and define
$k_i=p_i^2-Dq_i^2$. By \cref{lem:irrational}, $k_i\ne0$. Moreover,
\begin{equation}\label{eq:smallnorm}
 \abs{k_i}=\abs{p_i-q_is}(p_i+q_is)
 <\frac{Cq_i}{Q_i}\le C.
\end{equation}
Here $p_i\le Dq_i$ and $s<D$ give $p_i+q_is<Cq_i$.
The weaker integral bound $1\le\abs{k_i}\le C$ is sufficient below.

The denominators are separated by the known scales. If
$q_{i+1}\le Q_i$, then \eqref{eq:smallnorm} at $i+1$ would give
\[
 1\le\abs{k_{i+1}}<\frac{CQ_i}{CQ_i}=1.
\]
Thus
\begin{equation}\label{eq:separated}
 q_{i+1}>Q_i\ge q_i.
\end{equation}
This replaces an assertion about infinitely many distinct approximations by
a specific finite list with an explicit growth schedule.

For $m_i=\abs{k_i}$, let $\epsilon_i\in\{0,1\}$ record the sign of $k_i$.
Encode its signature by
\begin{equation}\label{eq:signature}
 \epsilon_iC^3+(m_i-1)C^2
 +(p_i\bmod m_i)C+(q_i\bmod m_i).
\end{equation}
Each of the last three digits is less than $C$, so the code is less than
$2C^3=M$. A second application of \cref{lem:pigeonhole} produces $i<j$
with equal signatures. Radix uniqueness, proved by division and remainder,
then gives
\[
 k_i=k_j=k\ne0,\qquad
 p_i\equiv p_j\pmod{\abs k},\qquad
 q_i\equiv q_j\pmod{\abs k}.
\]

\subsection{The quotient is an integral nontrivial unit}

Define
\[
 U=p_ip_j-Dq_iq_j,\qquad V=q_ip_j-p_iq_j.
\]
The congruences imply $\abs k\mid U,V$, since modulo $\abs k$ the first
quantity is $p_j^2-Dq_j^2=k$ and the second is zero. Thus $U/k,V/k$ are
integers. Expanding the products gives Brahmagupta's norm identity,
\begin{equation}\label{eq:quotientnorm}
 U^2-DV^2=(p_i^2-Dq_i^2)(p_j^2-Dq_j^2)=k^2.
\end{equation}
Consequently $(U/k)^2-D(V/k)^2=1$.

If $V=0$, then $p_iq_j=p_jq_i$. Squaring and substituting the two
equal-norm equations yields $k(q_j^2-q_i^2)=0$. Since $k\ne0$ and the
denominators are positive, $q_i=q_j$, contradicting \eqref{eq:separated}.
Therefore
\[
 x=\abs{U/k},\qquad y=\abs{V/k}
\]
satisfy $x^2-Dy^2=1$ and $y\ge1$; hence $x\ge2$.

For an explicit coordinate bound put $Q=C^{M+1}$. Every sample has
$q_i\le Q$ and $p_i\le DQ$. Even using the looser bound
$p_i\le(D+1)Q$ gives
\[
 \abs U\le(D^2+3D+1)Q^2,
 \qquad \abs V\le2(D+1)Q^2.
\]
Since $\abs k\ge1$, both output coordinates are at most
\begin{equation}\label{eq:H}
 \boxed{\collB(D)=(D+2)^2(2D+1)^{\,2(2(2D+1)^3+1)}.}
\end{equation}
This completes a uniform primitive recursive existence proof independently
of any continued-fraction theory. It also supplies a simple alternative
definition of the fundamental solution: scan $1\le y\le\collB(D)$ for
the first square $1+Dy^2$. The next section proves a much better algorithm.

\section{A finite continued-fraction theory}\label{sec:cf}

We construct finite integer recurrences and prove their properties. An
infinite continued fraction and a theorem about its convergence are not
prerequisites. Put $a=\isqrt(D)$ and retain the coded positive root $s$.

\subsection{Reduced states and a reversible transition}

\begin{definition}
A reduced state is a pair $(m,d)$ of integers with
\begin{equation}\label{eq:reduced}
 0\le m\le a,\quad 1\le d\le2a,\quad d\mid D-m^2,
 \qquad s-m<d<s+m.
\end{equation}
Its digit and successor are
\begin{equation}\label{eq:cfstep}
 c=\left\lfloor\frac{a+m}{d}\right\rfloor,
 \qquad m'=dc-m,
 \qquad d'=\frac{D-m'^2}{d}.
\end{equation}
The displayed floor is ordinary integer division.
\end{definition}

For the connection with ordinary continued-fraction notation, put
$\alpha_n=(s+m_n)/d_n$. Rationalizing the denominator in
\eqref{eq:cfstep} gives
\[
 \alpha_{n+1}=\frac1{\alpha_n-c_n}.
\]
Starting with $\alpha_0=s$, finite substitution therefore yields
$s=[c_0;c_1,\ldots,c_n,\alpha_{n+1}]$. This is an identity of finitely
many coded field operations; it requires no infinite-limit theorem.

The inequalities in \eqref{eq:reduced} are decidable arithmetic conditions.
For example, since $d+m>0$, $s-m<d$ means $D<(d+m)^2$;
$d<s+m$ is automatic for $d\le m$ and otherwise means
$(d-m)^2<D$. Thus the reduced-state test can be implemented without field
objects or floating-point calculations.

\begin{lemma}[Closure and injectivity]\label{lem:reduced}
The transition \eqref{eq:cfstep} maps reduced states to reduced states and
is injective on them.
\end{lemma}
\begin{proof}
Writing $a+m=dc+t$ with $0\le t<d$ and using $0<s-a<1$ shows
\[
 c=\left\lfloor\frac{s+m}{d}\right\rfloor.
\]
The upper reduced inequality gives $c\ge1$, and the floor inequalities give
\begin{equation}\label{eq:mprime}
 s-d<m'<s.
\end{equation}
Also $m'\ge0$. If $d\ge m$, use $dc-m\ge d-m$; if $d<m$, then
$d<m<s$ and $dc>s+m-d>m$.
Thus $0\le m'\le a$. Since $m'\equiv-m\pmod d$, $d$ divides
$D-m'^2$, which is positive. Therefore $d'$ is a positive integer.

We have $m'\ge d-m$ because $c\ge1$, so $s+m'>d$; and
\eqref{eq:mprime} says $s-m'<d$. In
\[
 d'=\frac{(s-m')(s+m')}{d}
\]
these two inequalities imply
$s-m'<d'<s+m'$. Finally $d'<s+a<2a+1$, so $d'\le2a$.
The relation $dd'=D-m'^2$ also gives the required divisibility by $d'$.
All reduced conditions are preserved.

For injectivity, a reduced predecessor $(m,d)$ of $(m',d')$ must have
\[
 d=(D-m'^2)/d'.
\]
The old inequalities imply $s-d<m<s$. With $m=dc-m'$ this gives
\[
 \frac{s+m'}d-1<c<\frac{s+m'}d.
\]
The upper endpoint is not an integer by \cref{lem:irrational}; hence
$c=\lfloor(s+m')/d\rfloor=\lfloor(a+m')/d\rfloor$. This recovers $c$
and then $m$ uniquely. Thus two reduced predecessors cannot differ.
\end{proof}

\subsection{The period has an explicit bound}

Start with
\[
 (m_0,d_0)=(0,1),\quad c_0=a,
 \qquad (m_1,d_1)=(a,D-a^2).
\]
The state at index $1$ is reduced: $D-a^2=(s-a)(s+a)$ and
$0<s-a<1<s+a$ imply $s-a<d_1<s+a$, and $d_1\le2a$ follows from
$a^2<D<(a+1)^2$.
There are at most
\begin{equation}\label{eq:S}
 S=2a(a+1)
\end{equation}
possible reduced states, counted by the rectangle in \eqref{eq:reduced}.
Construct exactly $S+1$ states starting with state $1$. The finite
pigeonhole lemma gives two equal states at indices $i<j$. Injectivity lets
us cancel $i-1$ transitions and obtain equality of state $1$ and state
$j-i+1$. Therefore a positive return time at most $S$ exists.

Let $\ell$ be the first return time, found by scanning $1,\ldots,S$.
All subsequent state equalities of period $\ell$ follow by induction on
the explicit recurrence. This is a finite induction at any requested
index, not an appeal to an infinite combinatorial principle.
The state $(a,1)$ is reduced, has digit $2a$, and maps to
$(a,D-a^2)$. By uniqueness of reduced predecessors,
\begin{equation}\label{eq:endpoint}
 (m_\ell,d_\ell)=(a,1),\qquad c_\ell=2a.
\end{equation}
The digit $2a$ occurs only at this state: $c=2a$ implies
$2ad\le a+m\le2a$, whence $d=1$ and $m=a$.
Thus the least state period is also the least repeating digit-block length.
Indeed, a shorter digit period would repeat the digit at index $\ell$ at
an earlier positive index, contradicting the first return.

More precisely, for every $k\ge1$,
\begin{equation}\label{eq:unitdenom}
 d_k=1\quad\Longleftrightarrow\quad\ell\mid k.
\end{equation}
Indeed, reducedness with $d_k=1$ gives $s-1<m_k<s$, so $m_k=a$ and
the next state is state $1$. Conversely a return at state $k+1$ has
predecessor $(a,1)$. Divide $k$ by $\ell$; recurrence periodicity reduces
the return to its remainder, which is zero by minimality of $\ell$.
This proves \eqref{eq:unitdenom} with finite division and finite induction.

\subsection{Convergents and an exact norm identity}

Define integers by
\begin{equation}\label{eq:convergents}
\begin{aligned}
 p_{-2}=0,&\quad p_{-1}=1,& p_n&=c_np_{n-1}+p_{n-2},\\
 q_{-2}=1,&\quad q_{-1}=0,& q_n&=c_nq_{n-1}+q_{n-2}.
\end{aligned}
\end{equation}
Negative subscripts name the four initial values; an implementation stores
them as two registers, so no negative recursion parameter is involved.
For $n\ge0$, $p_n,q_n$ are positive. The $p_n$ strictly increase,
$q_0=1$, $q_1=c_1\ge1$, and $q_{n+1}>q_n$ for $n\ge1$.
In particular $q_n\ge n$ for $n\ge1$.

\begin{lemma}[Simultaneous integer invariants]\label{lem:cfinvariants}
For every $n\ge0$,
\begin{align}
 p_{n-1}&=m_nq_{n-1}+d_nq_{n-2},\label{eq:inv1}\\
 Dq_{n-1}&=m_np_{n-1}+d_np_{n-2},\label{eq:inv2}\\
 p_{n-1}q_{n-2}-p_{n-2}q_{n-1}&=(-1)^n.\label{eq:det}
\end{align}
Consequently
\begin{equation}\label{eq:cfnorm}
 \boxed{p_{n-1}^2-Dq_{n-1}^2=(-1)^n d_n.}
\end{equation}
\end{lemma}
\begin{proof}
At $n=0$ substitute the four initial values and $(m_0,d_0)=(0,1)$.
For the induction step write $(m,d,c)=(m_n,d_n,c_n)$ and
$m'=dc-m$, $d'=(D-m'^2)/d$. The old identities and the recurrences give
\[
 dp_n=m'p_{n-1}+Dq_{n-1},\qquad
 dq_n=p_{n-1}+m'q_{n-1}.
\]
Subtract $m'$ times the second equation from the first:
\[
 d(p_n-m'q_n)=(D-m'^2)q_{n-1}=dd'q_{n-1}.
\]
Cancel the positive $d$ to obtain \eqref{eq:inv1} at the next index.
Similarly,
\[
 d(Dq_n-m'p_n)=(D-m'^2)p_{n-1}=dd'p_{n-1}
\]
gives \eqref{eq:inv2}. The determinant changes sign under
\eqref{eq:convergents}, proving \eqref{eq:det}.
Finally eliminate $m_n$ between \eqref{eq:inv1} and \eqref{eq:inv2}:
\[
 p_{n-1}^2-Dq_{n-1}^2
 =d_n(p_{n-1}q_{n-2}-p_{n-2}q_{n-1}).
\]
This is \eqref{eq:cfnorm}. The whole simultaneous invariant consists of
three polynomial equalities.
\end{proof}

At $n=\ell$, \eqref{eq:cfnorm} has right side $(-1)^\ell$.
If $\ell$ is even, it already supplies a positive norm-one solution.
If $\ell$ is odd, use $n=2\ell$, for which $d_{2\ell}=1$.
Thus at most $2S$ convergents suffice for existence.

All digits are at most $2a$. The recurrence gives, by elementary induction,
\begin{equation}\label{eq:convbound}
 p_n,q_n\le(2a+1)^{n+1}\qquad(n\ge0).
\end{equation}
For the induction step, with $K=2a+1$, use
$2aK^{n+1}+K^n\le K^{n+2}$; the initial two registers are checked directly.
We obtain the smaller explicit bound
\begin{equation}\label{eq:B}
 \boxed{\CFB(D)=(2\isqrt(D)+1)^{4\isqrt(D)(\isqrt(D)+1)}.}
\end{equation}

\subsection{A finite proof of Legendre's criterion}\label{sec:legendre}

To establish minimality and the converse of the negative-Pell criterion,
we need to show that an arbitrary sufficiently accurate rational
approximation occurs in the recurrence. This can also be done finitely.
Put $\varepsilon_n=q_ns-p_n$. The norm identity shows that consecutive
$\varepsilon_n$ have opposite signs: since $p_n+q_ns>0$, the sign of
$p_n-q_ns$ is $(-1)^{n+1}$.

\begin{lemma}[Finite best approximation of the second kind]\label{lem:bestsecond}
Let $q\ge1$ and choose $n$ so $q_n\le q<q_{n+1}$.
For every integer $p$,
\[
 \abs{qs-p}\ge\abs{q_ns-p_n}.
\]
Such an $n$ can be found by examining the indices $0,\ldots,q+1$.
\end{lemma}
\begin{proof}
We have $q_0=1\le q$ and $q_{q+1}\ge q+1>q$, so take the last index
with $q_n\le q$. This handles the possible equality $q_0=q_1$.
The determinant of $(p_n,q_n)$ and $(p_{n+1},q_{n+1})$ is $\pm1$.
Its inverse integer matrix therefore gives integers $h,k$ with
\[
 (p,q)=h(p_n,q_n)+k(p_{n+1},q_{n+1}).
\]
If $k=0$, then $h\ge1$, and the required error comparison is immediate.
If $k\ge1$, the inequality $q<q_{n+1}$ forces $h<0$.
If $k\le-1$, positivity of $q$ forces $h>0$.
Thus for $k\ne0$ the coefficients have opposite signs, with $\abs h\ge1$.
Since the two convergent errors also have opposite signs, the summands in
\[
 qs-p=h\varepsilon_n+k\varepsilon_{n+1}
\]
have the same sign. Its absolute value is therefore at least
$\abs h\abs{\varepsilon_n}\ge\abs{\varepsilon_n}$.
\end{proof}

\begin{theorem}[Legendre, in finite form]\label{thm:legendre}
If $p\in\Z$, $q\ge1$, and
\[
 \abs{s-p/q}<\frac1{2q^2},
\]
then $p/q=p_n/q_n$ for an index $n\le q$.
\end{theorem}
\begin{proof}
Choose $n$ as in \cref{lem:bestsecond} and put $E=\abs{qs-p}$.
If the two rational numbers differed, their nonzero integer determinant
would satisfy
\[
\begin{aligned}
 1&\le\abs{pq_n-qp_n}\\
  &\le q_n\abs{p-qs}+q\abs{q_ns-p_n}\\
  &\le(q_n+q)E\le2qE<1,
\end{aligned}
\]
which is impossible. The test whether the determinant is zero is decidable,
so the argument yields its equality to zero intuitionistically.
\end{proof}

\subsection{Fundamental solutions and period parity}\label{sec:cfclassification}

Suppose $p,q>0$ and $p^2-Dq^2=\pm1$. A common divisor of $p,q$ has
square dividing $1$, so the pair is coprime. Also
$p^2\ge Dq^2-1\ge q^2$, hence $p\ge q$. Therefore
\[
 \abs{s-p/q}=\frac1{q(p+qs)}<\frac1{2q^2}
\]
because $s>1$. By \cref{thm:legendre}, $p/q$ is a convergent.
The convergent pair is coprime by its determinant identity, so equality of
the positive reduced fractions implies $(p,q)=(p_n,q_n)$ exactly.

By \eqref{eq:cfnorm} and \eqref{eq:unitdenom}, the convergents of norm
$\pm1$ are precisely those at indices $k-1$ with $k$ a positive multiple
of $\ell$, and their norm is $(-1)^k$. We have proved both directions of
the classical criterion:
\begin{theorem}\label{thm:parity}
Negative Pell is solvable if and only if $\ell$ is odd. The fundamental
positive norm-one solution is
\begin{equation}\label{eq:fundCF}
 (A,B)=
 \begin{cases}
 (p_{\ell-1},q_{\ell-1}),&\ell\text{ even},\\
 (p_{2\ell-1},q_{2\ell-1}),&\ell\text{ odd}.
 \end{cases}
\end{equation}
When $\ell$ is odd, the least positive negative-Pell solution is
$(p_{\ell-1},q_{\ell-1})$.
\end{theorem}
\begin{proof}
If $\ell$ is even, all its multiples are even, so no convergent has norm
$-1$; the preceding paragraph excludes every other positive solution too.
If $\ell$ is odd, $k=\ell$ supplies a negative solution, and the first
even multiple is $2\ell$. In each case \eqref{eq:fundCF} is the first
norm-one convergent. The numerators $p_n$ strictly increase. Within either
fixed norm, $x^2=Dy^2\pm1$ shows that ordering positive $x$ is equivalent
to ordering positive $y$. Thus the first eligible convergent is fundamental.
For integer solutions, absolute values reduce to the positive case.
\end{proof}

This proof supplies a total bounded numerical algorithm, the exact
period-parity theorem, and a uniform size bound. It does not appeal to
Lagrange's periodicity theorem, Legendre's theorem, or a general theorem
about irrational reals as an external premise; the needed finite versions
have all been proved above.

\section{Fundamental units and every norm-one solution}\label{sec:units}

\subsection{Pair multiplication and the power sequence}

For integer pairs write
\[
 (x,y)(r,t)=(xr+Dyt,xt+yr),\qquad
 \overline{(x,y)}=(x,-y),\qquad
 \Nm(x,y)=x^2-Dy^2.
\]
Expansion proves commutativity, associativity, identity $(1,0)$, and
\begin{equation}\label{eq:normmul}
 \Nm(zw)=\Nm(z)\Nm(w),\qquad z\bar z=(\Nm(z),0).
\end{equation}
Thus a norm-one pair has inverse its conjugate. This is the same arithmetic
as the rational-pair field, restricted to integer coordinates.

Let $u=(A,B)$ be the fundamental solution from \cref{thm:parity}.
Alternatively choose the least successful $y$ in the bounded search
$1\le y\le\collB(D)$ from \cref{sec:collision}; the definition is the
same by minimality. In either construction
\[
 A\ge2,\qquad B\ge1,\qquad B<A,\qquad A^2-DB^2=1.
\]
Define $u^n=(X_n,Y_n)$ by the two recurrences in \cref{thm:main}.
Quantifier-free induction using \eqref{eq:normmul} gives
\begin{equation}\label{eq:powerfacts}
 X_n^2-DY_n^2=1,\qquad X_n\ge1,\qquad Y_n\ge0.
\end{equation}
Furthermore
\[
 Y_{n+1}=BX_n+AY_n\ge2Y_n+1,
\]
so $Y_n\ge2^n-1\ge n$ and the $Y_n$ strictly increase.
For an upper growth bound put $G=A+DB$. Both coefficients in
\[
 X_{n+1}+Y_{n+1}=(A+B)X_n+(DB+A)Y_n
\]
are at most $G$, hence $X_n+Y_n\le G^n$.

For fixed $m$, induction on $n$ or pair associativity proves
\begin{equation}\label{eq:addition}
\begin{aligned}
 X_{m+n}&=X_mX_n+DY_mY_n,\\
 Y_{m+n}&=X_mY_n+Y_mX_n.
\end{aligned}
\end{equation}
All occurrences of powers below can be read as these explicit recurrences.

\subsection{A fixed-fuel descent proves completeness}

\begin{lemma}[Positive unit descent]\label{lem:unitdescent}
If $x,y\ge0$, $x^2-Dy^2=1$, and $y>0$, then
\[
 r=Ax-DBy,\qquad t=Ay-Bx
\]
satisfy $r>0$, $0\le t<y$, $r^2-Dt^2=1$, and
$(x,y)=u(r,t)$.
\end{lemma}
\begin{proof}
Fundamental minimality gives $y\ge B$. Compare squares of the indicated
nonnegative quantities:
\begin{align*}
 (Ax)^2-(DBy)^2&=Dy^2+A^2>0,\\
 (Ay)^2-(Bx)^2&=y^2-B^2\ge0,\\
 (Bx)^2-((A-1)y)^2&=B^2+2(A-1)y^2>0.
\end{align*}
The first line gives $r>0$, the second $t\ge0$, and the third
$Bx>(A-1)y$, or $t<y$. Norm preservation and the reconstruction equation
follow by multiplication by $u^{-1}=(A,-B)$ and then by $u$.
\end{proof}

\begin{theorem}[All solutions, with a bounded exponent algorithm]\label{thm:powers}
For every nonnegative norm-one solution $(x,y)$ there is a unique $n\le y$
with $(x,y)=(X_n,Y_n)$. The exponent is primitive recursively computable
from $D,x,y$.
\end{theorem}
\begin{proof}
Run exactly the input $y$ stages. The state stores a residual pair and a
counter $c$, initially $(x,y)$ and $0$. When the residual second
coordinate is positive, apply \cref{lem:unitdescent} and increment $c$;
when it is zero, leave the whole state fixed.
After $j$ stages the quantifier-free invariant is
\[
 (x,y)=u^c(\text{residual pair}),\quad
 \Nm(\text{residual pair})=1,
\]
together with nonnegative residual coordinates, $c\le j$, and residual
second coordinate at most $\max(y-j,0)$.
The descent lemma proves its preservation. At stage $y$ the residual second
coordinate is zero, so its first coordinate is $1$. Thus the stored counter
is the desired $n\le y$. Strict growth of $Y_n$ proves uniqueness.
\end{proof}

In particular, ``infinitely many solutions'' has the stronger constructive
meaning of an explicit injective primitive recursive sequence of solutions.
For arbitrary integer solutions, take $(\abs x,\abs y)$ and then restore
the independent signs. Equivalently the norm-one group is
\[
 \{\,\pm u^k:k\in\Z\,\},
\]
where $u^{-n}=(X_n,-Y_n)$. This set notation abbreviates the coordinate
classification, rather than an additional set-existence assumption in iPRA.

\section{Negative Pell and square roots of the fundamental unit}\label{sec:negative}

The period criterion already decides the equation. A second proof by
integer descent gives a useful bounded test and explains its structure.

\begin{lemma}[Negative descent]\label{lem:negdescent}
Suppose $x,y\ge0$, $x^2-Dy^2=-1$, and $y>B$. Then
$(r,t)=(Ax-DBy,Ay-Bx)$ has positive coordinates, norm $-1$, and $t<y$.
\end{lemma}
\begin{proof}
First $x,y\ge1$, since $D\ge2$. The square comparisons now are
\begin{align*}
 (Ax)^2-(DBy)^2&=D(y^2-B^2)-1>0,\\
 (Ay)^2-(Bx)^2&=y^2+B^2>0,\\
 (Bx)^2-((A-1)y)^2&=2(A-1)y^2-B^2>0.
\end{align*}
The first inequality uses $y\ge B+1$ and $D\ge2$; the third follows from
$A\ge2$ and $y>B$. The same comparisons as before give $r,t>0$ and
$t<y$. Norm preservation follows from \eqref{eq:normmul}.
\end{proof}

The descent with the original $y$ as fuel reaches a positive solution
$(R,T)$ with $1\le T\le B$. Conversely such a pair is itself a solution.
Thus
\begin{equation}\label{eq:negativefinite}
 x^2-Dy^2=-1\text{ is solvable}
 \quad\Longleftrightarrow\quad
 \exists\,1\le T\le B:\ DT^2-1\text{ is a square}.
\end{equation}
This is a bounded decision rule; no undecided negative case is left to an
unbounded search.

\begin{theorem}[Negative units]\label{thm:negativeunits}
There is at most one positive solution $v=(R,T)$ with $T\le B$.
If it exists, $v^2=u$, and all positive negative-Pell solutions are
$vu^n=v^{2n+1}$ for $n\ge0$.
\end{theorem}
\begin{proof}
For such a solution,
$R^2=DT^2-1\le DB^2-1=A^2-2$, so $R<A$.
Its square is a positive norm-one solution. By \cref{thm:powers},
$v^2=u^k$ for some $k\ge1$. Its second coordinate satisfies
\[
 0<2RT<2AB=Y_2.
\]
Strict growth of $Y_k$ forces $k=1$. In coordinates this gives
\begin{equation}\label{eq:negativesquare}
 A=2R^2+1,\qquad B=2RT,\qquad DT^2=R^2+1.
\end{equation}
These equations determine $R>0$ uniquely from $A$, and then $T>0$
uniquely from $B$, proving uniqueness of the bounded seed.
Reverse the fixed-fuel descent in \cref{lem:negdescent}: every positive
negative solution is $vu^n$. Conversely each such product has positive
coordinates and norm $-1$. Since $u=v^2$, it is also $v^{2n+1}$.
\end{proof}

Equations \eqref{eq:negativesquare} give a further finite test: check for
positive integers $R,T$ satisfying all three equalities. Testing only that
$(A-1)/2$ is an integer square is insufficient for general nonsquare $D$.
For example $D=20$ has fundamental solution $(A,B)=(9,2)$ and
$(A-1)/2=4$, but $20T^2=5$ has no integer solution. This example also
shows why keeping the actual order $\Z[\sqrt D]$ is important.

For cheap preliminary obstructions, $D\equiv0$ or $3\pmod4$ rules out
negative Pell by the two square residues modulo $4$. These tests are
optional optimizations; the period test and \eqref{eq:negativefinite}
are complete for every allowed $D$.

\section{Every generalized Pell equation has a finite decision procedure}\label{sec:general}

Fix a nonzero integer $N$ and put
\begin{equation}\label{eq:seedbounds}
 K=\abs N\ge1,\qquad T=AK,\qquad R=(DA+1)K=DT+K.
\end{equation}
Here $A,B$ are the fundamental norm-one coordinates. In fact any explicitly
given positive norm-one unit would suffice for the following reduction;
using the fundamental one usually makes the seed box smaller.

\begin{lemma}[Uniform reduction for either sign of $N$]\label{lem:generaldescent}
Suppose $x,y\ge0$, $x^2-Dy^2=N$, and $y>T$. Then multiplication by
$(A,-B)$ produces a pair with positive coordinates, norm $N$, and second
coordinate strictly less than $y$.
\end{lemma}
\begin{proof}
The three square differences are
\begin{align*}
 (Ay)^2-(Bx)^2&=y^2-B^2N\ge y^2-B^2K>0,\\
 (Ax)^2-(DBy)^2&=Dy^2+A^2N\ge Dy^2-A^2K>0,\\
 (Bx)^2-((A-1)y)^2
   &=2(A-1)y^2+B^2N\ge2(A-1)y^2-B^2K>0.
\end{align*}
For the strict inequalities use $y>AK$, $A>B$, $K\ge1$, and $A\ge2$.
In particular $y^2>A^2K^2\ge A^2K>B^2K$. The comparisons give
$Ay-Bx>0$, $Ax-DBy>0$, and $Ay-Bx<y$, respectively.
The norm is preserved by \eqref{eq:normmul}.
\end{proof}

Apply this step while the second coordinate exceeds $T$, with exactly the
original $y$ stages of fuel and a stationary final state. The terminal pair
$(r,t)$ satisfies $0\le t\le T$ and
\[
 r^2=Dt^2+N\le DT^2+K\le(DT+K)^2=R^2.
\]
The last inequality follows from
\[
 (DT+K)^2-(DT^2+K)
 =D(D-1)T^2+2DTK+K(K-1)\ge0.
\]
Thus $0\le r\le R$.

Define the finite seed list
\begin{equation}\label{eq:seedset}
 \mathcal S_{D,N}=\{(r,t):0\le r\le R,\ 0\le t\le T,
                         \ r^2-Dt^2=N\}.
\end{equation}
This notation means a bounded scan and a coded output list. One may scan
just $t=0,\ldots,T$: if $Dt^2+N$ is nonnegative, compute its integer
square root and test its square. The bound on $r$ has already been proved.

\begin{theorem}[Finite seeds and uniform decision]\label{thm:general}
For $N\ne0$, integer solvability of $x^2-Dy^2=N$ is equivalent to
nonemptiness of \eqref{eq:seedset}. Every nonnegative solution has a
representation
\[
 (x,y)=(r,t)u^n,\qquad (r,t)\in\mathcal S_{D,N},\quad n\le y,
\]
computed by fixed-fuel descent. Conversely every such product is a
nonnegative solution. Independent signs on the two coordinates give every
integer solution.
\end{theorem}
\begin{proof}
Taking absolute values preserves $x^2-Dy^2$, so consider nonnegative inputs.
The preceding descent and a counter yield a seed. The reconstruction
invariant is the same as in \cref{thm:powers}, with the threshold $T$ in
place of zero: the input equals $u^c$ times the current pair.
If the state were still active after $y$ stages, its positive second
coordinate would have decreased at least $y$ times from $y$, which is
impossible. Equivalently the positive excess above $T$ decreases at every
active stage and is zero at the final stage. Thus the seed is in the stated
box and reversing the steps proves the representation.
Conversely the seed norm is $N$, the power norm is one, and their
nonnegative coordinates multiply to nonnegative coordinates. This proves
all claims.
\end{proof}

The finite list need not have one seed per orbit: distinct seeds can belong
to the same unit orbit. The theorem claims a finite, complete covering by
explicitly generated families. This suffices for decision and exhaustive
enumeration. For the full signed unit action, adjoin signs and conjugates
to the list if a set of covering representatives is desired.

When $N=0$, \cref{lem:irrational} excludes any $y\ne0$; then $x=0$.
Thus the decision algorithm treats zero separately and terminates for every
integer $N$. The explicit function $A(D)$ makes the bounds in
\eqref{eq:seedbounds} primitive recursive jointly in $D$ and a signed code
for $N$. This is stronger than a procedure that merely keeps searching until
a solution happens to appear.

\subsection{Degenerate discriminants}

The nonsquare guard is decidable by $a=\isqrt(D)$. For positive square
$D=c^2$, the ordinary equation factors as
$(x-cy)(x+cy)=1$, giving only $(x,y)=(\pm1,0)$.
For $D=0$, the ordinary equation has $x=\pm1$ and arbitrary $y$.
Our algorithms return a designated invalid-input result on these excluded
discriminants; they do not execute a search justified only for nonsquares.
The main generalized theorem concerns \eqref{eq:guard}, so no exceptional
square-discriminant factorization cases are hidden in its statement.

\section{Recurrences, divisibility, and finite congruence periods}\label{sec:arithmeticstructure}

\subsection{Useful identities}

The addition formulas and $A^2-DB^2=1$ give
\begin{align}
 X_{n+2}+X_n&=2AX_{n+1},&
 Y_{n+2}+Y_n&=2AY_{n+1},\label{eq:secondorder}\\
 X_{2n}&=2X_n^2-1=1+2DY_n^2,&
 Y_{2n}&=2X_nY_n.\label{eq:doubling}
\end{align}
For example, substituting the first-order recurrences twice makes the
coefficient of $X_n$ in $X_{n+2}$ equal to
$A^2+DB^2=2A^2-1$, and that of $Y_n$ equal to $2ADB$;
this is $2AX_{n+1}-X_n$. The second-coordinate calculation is identical
in form. The doubling identities follow directly from \eqref{eq:addition}.
Writing the second-order recurrences additively avoids truncated subtraction
when they are encoded as equations of natural-number terms.

Another useful determinant identity is
\begin{equation}\label{eq:powerdet}
 X_mY_{m+k}-X_{m+k}Y_m=Y_k.
\end{equation}
Substitute \eqref{eq:addition}; the remaining coefficient is
$X_m^2-DY_m^2=1$. In particular the determinant between consecutive
power pairs is $B$.

\subsection{Strong divisibility}

\begin{theorem}\label{thm:gcd}
For all $m,n\ge0$,
\[
 \gcd(Y_m,Y_n)=Y_{\gcd(m,n)}.
\]
For $m\ge1$ and $n\ge0$,
\[
 Y_m\mid Y_n\quad\Longleftrightarrow\quad m\mid n.
\]
\end{theorem}
\begin{proof}
A common divisor of $X_m,Y_m$ has square dividing
$X_m^2-DY_m^2=1$, so $\gcd(X_m,Y_m)=1$. The addition formula gives
\[
\begin{aligned}
 \gcd(Y_m,Y_{m+k})
 &=\gcd(Y_m,X_mY_k+X_kY_m)\\
 &=\gcd(Y_m,X_mY_k)\\
 &=\gcd(Y_m,Y_k).
\end{aligned}
\]
The last cancellation follows from a Bezout identity for $X_m,Y_m$:
any divisor of $Y_m$ and $X_mY_k$ also divides $Y_k$, and the converse
divisibility is immediate.

Run the subtraction Euclidean algorithm on $(m,n)$, swapping as needed,
with at most $m+n$ nonterminal steps. The displayed equality preserves
the gcd of the corresponding $Y$ coordinates at every subtraction.
At the terminal index pair $(g,0)$ or $(0,g)$, where $g=\gcd(m,n)$,
the gcd is $Y_g$, since $Y_0=0$. If $m=n=0$, both sides are zero.
This proves the first assertion by a bounded loop invariant.

If $m\mid n$, then $\gcd(m,n)=m$, so $Y_m\mid Y_n$.
Conversely, because $m\ge1$ and $Y_m>0$, divisibility gives
$Y_m=\gcd(Y_m,Y_n)=Y_{\gcd(m,n)}$. Strict growth of the sequence implies
$m=\gcd(m,n)$, and hence $m\mid n$.
\end{proof}

The second coordinates used here are the actual Pell coordinates, without
normalizing by $B$. Since $Y_1=B$, the theorem in particular proves
$B\mid Y_n$ for every $n$.

\subsection{A finite replacement for infinite-domain pigeonhole}

\begin{theorem}\label{thm:modular}
For every $h\ge1$ one can compute $t$ with $1\le t\le h^2$ such that
\[
 X_t\equiv1\pmod h,\qquad Y_t\equiv0\pmod h.
\]
Consequently $(X_{n+t},Y_{n+t})\equiv(X_n,Y_n)\pmod h$ for every $n$.
\end{theorem}
\begin{proof}
Encode the residue pair at index $j$ by
$h(X_j\bmod h)+(Y_j\bmod h)<h^2$.
Among the first $h^2+1$ indices $0,\ldots,h^2$, the finite pigeonhole
lemma gives equal states at $i<j$. Multiplication by $u$ on residue pairs
is invertible: its matrix is
\[
 \begin{pmatrix}A&DB\\B&A\end{pmatrix},\qquad
 \det=A^2-DB^2=1,
\]
and multiplication by $(A,-B)$ is its inverse, also modulo $h$.
Cancel $i$ multiplications to get the identity state at $t=j-i$.
The addition formula then proves the asserted period at every index.

For a direct implementation, simply scan $1,\ldots,h^2$ for the first
identity state. The collision proof guarantees success of that bounded
scan. It includes $h=1$: every residue is zero and $t=1$ is valid.
\end{proof}

This is the appropriate arithmetic substitute for a theorem that a map from
an infinite domain to a finite set repeats. It returns an explicit finite
period bound and a specific exponent. In particular it proves, with a
primitive recursive witness, that some positive $Y_t$ is divisible by any
given positive modulus.

\section{Rational approximation with a primitive recursive modulus}\label{sec:approximation}

\subsection{Nested enclosures}

For $n\ge1$ define positive rationals
\[
 L_n=\frac{DY_n}{X_n},\qquad U_n=\frac{X_n}{Y_n}.
\]
The norm identity proves
\begin{equation}\label{eq:enclosure}
 L_n^2<D<U_n^2,\qquad U_n-L_n=\frac1{X_nY_n}.
\end{equation}
These are statements about integer cross products. In particular they have
meaning before constructing any notion of a completed real line.
Using \eqref{eq:powerdet} with $k=1$ gives
\begin{equation}\label{eq:nesting}
 U_n-U_{n+1}=\frac B{Y_nY_{n+1}}>0,\qquad
 L_{n+1}-L_n=\frac{DB}{X_nX_{n+1}}>0.
\end{equation}
The intervals are therefore nested. Since $X_n\ge Y_n$ and
$Y_n\ge2^n-1\ge2^{n-1}$,
\begin{equation}\label{eq:modulus}
 U_n-L_n\le2^{-2n+2}.
\end{equation}
For accuracy $2^{-k}$, the explicit choice $n=k+1$ suffices, and even
gives width at most $2^{-2k}$.

If desired, the sequence $U_{k+1}$ is a constructive Cauchy presentation of
the positive root. For $j\ge k$, nesting bounds the difference between
its $j$th and $k$th values by $2^{-2k}$. Its squares approach $D$ with
an explicit modulus as well: from $U_n^2-D=1/Y_n^2$ and the same growth
bound, the error is at most $2^{-2n+2}$. Internally, the theorem is this
family of rational inequalities and the named modulus. No principle
asserting limits for arbitrary Cauchy sequences is used.

\subsection{One-sided best approximation}

\begin{theorem}\label{thm:onesided}
Let $n\ge1$. If $p\ge0$, $1\le q\le Y_n$, and $p^2>Dq^2$, then
\[
 \frac pq\ge\frac{X_n}{Y_n},
\]
with strict inequality if $q<Y_n$. If $R,T>0$ satisfy
$R^2-DT^2=-1$, then for $p\ge0$, $1\le q\le T$, and $p^2<Dq^2$,
\[
 \frac pq\le\frac RT,
\]
again strictly when $q<T$.
\end{theorem}
\begin{proof}
In the upper case $p^2-Dq^2$ is an integer at least $1$, so
\[
 (pY_n)^2-(X_nq)^2
 =(p^2-Dq^2)Y_n^2-q^2\ge Y_n^2-q^2\ge0.
\]
All factors are nonnegative; comparison of squares gives the assertion,
and $q<Y_n$ makes it strict.
In the lower case $Dq^2-p^2\ge1$, and
\[
 (Rq)^2-(pT)^2=(Dq^2-p^2)T^2-q^2\ge T^2-q^2\ge0.
\]
The same square comparison proves the lower statement.
\end{proof}

These are one-sided optimality statements. They do not assert that a Pell
approximant is closest among rationals on both sides of the root. Their
precise inequalities are particularly simple to encode in the base theory.

\section{Examples and exact certificates}\label{sec:examples}

The reference implementation produces the following values using exact
integer arithmetic. The period is the first return of the reduced state.
The dash means that the complete parity test rules out negative Pell.

\begin{center}
\small
\begin{tabular}{rrrrr}
\toprule
$D$ & $\ell$ & $A$ & $B$ & Least positive norm $-1$ pair\\
\midrule
2 & 1 & 3 & 2 & $(1,1)$\\
3 & 2 & 2 & 1 & ---\\
7 & 4 & 8 & 3 & ---\\
12 & 2 & 7 & 2 & ---\\
13 & 5 & 649 & 180 & $(18,5)$\\
34 & 4 & 35 & 6 & ---\\
61 & 11 & 1\,766\,319\,049 & 226\,153\,980 & $(29\,718,3\,805)$\\
\bottomrule
\end{tabular}
\end{center}

For $D=13$, $a=3$, the reduced states during one period are
\[
 (3,4),\ (1,3),\ (2,3),\ (1,4),\ (3,1),
\]
with digits $1,1,1,1,6$. The initial digit is $3$.
The convergent before the terminal digit is $18/5$, with
$18^2-13\cdot5^2=-1$. Its pair square is
\[
 (18,5)^2=(18^2+13\cdot5^2,\,2\cdot18\cdot5)=(649,180),
\]
the fundamental positive unit. The recurrence begins
$(1,0),(649,180),\ldots$ and generates all positive norm-one solutions.

As a generalized example, $D=2$, $N=-7$, and $u=(3,2)$ give
$T=21$, $R=49$. The pairs $(1,2)$ and $(5,4)$ are seeds because their
norms are $-7$, and each generates a family by multiplication by $(3,2)$.
For example
\[
 (1,2)(3,2)=(11,8),\qquad 11^2-2\cdot8^2=-7.
\]
The complete seed scan may include further members of these same families;
it is intentionally a complete finite covering, not a minimal orbit list.

For $D=2$ and $n=1$, the enclosure is $[4/3,3/2]$, of width $1/6$.
At $n=2$, it is $[24/17,17/12]$, of width $1/204$.
Both are obtained by cross multiplication of Pell coordinates, with no
floating-point square root.

\section{Why the proofs fit iPRA}\label{sec:audit}

\subsection{Fixed programs and totalized operations}

The relevant function symbols form a fixed finite collection: signed
arithmetic, rational-pair arithmetic, pairing and finite arrays, division
with remainder, square root by bounded search, finite Boolean folds,
bounded collision selection, the two Pell algorithms, and the various
power and descent recurrences. Each is built from zero, successor,
projections, composition, and primitive recursion.

All functions must have values on all numerical inputs. Define division by
zero, unsuccessful selectors, and invalid-discriminant solver calls to
return designated default codes. The correctness theorem has the guard
\eqref{eq:guard}, and its invariants prove that the defaults for invalid
intermediate states are never used on guarded inputs. Python's validation
exceptions in the reference implementation serve as diagnostics; the
object-level functions use default natural values instead.

The term ``fixed'' is essential. We are not assuming a primitive recursive
universal evaluator for arbitrary primitive recursive program codes.
The metatheory constructs a particular syntax tree of primitive recursive
definitions for each program used here. Parameters $D,N,x,y,h,k$ are
ordinary numerical inputs to those fixed programs.

\subsection{Induction invariants and finite quantifiers}

The main induction instances can be arranged as follows. Every entry in
the third column is a quantifier-free combination of equations and
inequalities between primitive recursive terms. A bounded assertion about
a trace is replaced by its characteristic Boolean fold.

\begin{longtable}{p{0.23\textwidth}p{0.22\textwidth}p{0.46\textwidth}}
\toprule
Construction & Induction counter & Invariant and witness data\\
\midrule
\endhead
Integer square root & Trial index, at most input & Current accepted square,
maximal accepted trial, finite-prefix test.\\
Euclid and Bezout & At most input sum plus one & Preserved gcd condition,
two linear combinations, decrease of second coordinate while active.\\
Finite collision & At most list length & Occupancy bits and count, success
flag, stored first indices, equality at a returned pair.\\
Dirichlet bins & At most $Q+1$ & Exact floor-square inequalities, bin in
$[0,Q)$, computed collision and cross-multiplied error tests.\\
Norm signatures & At most $M+1$ & Nonzero bounded norm, denominator
separation, radix code below $M$, selected matching indices.\\
Reduced CF states & At most requested index & Conditions
\eqref{eq:reduced}; exact divisibility; predecessor recovery equation.\\
CF convergents & At most requested index & The three simultaneous
identities \eqref{eq:inv1}--\eqref{eq:det}, positivity, coordinate bounds.\\
Period cancellation & At most $S$ or a given index & Explicit equality of
two states moved backwards, or equality after a period shift.\\
Legendre bracketing & At most $q+1$ & Last denominator no larger than $q$;
determinant and error-sign identities.\\
Power sequence & Exponent & Norm one, growth, addition and determinant
identities with other indices as free parameters.\\
Unit descent & Original second coordinate & Preserved norm and signs,
strict decrease while active, input equals accumulated power times residual.\\
General-norm descent & Original second coordinate & The same reconstruction;
positive excess above $A\abs N$ decreases while active.\\
Strong divisibility & At most $m+n$ & Gcd of the two current Pell
coordinates unchanged during index Euclid.\\
Residue period & At most $h^2+1$ & Pair residue code, collision witnesses,
inverse multiplication and reconstruction of the difference exponent.\\
Rational enclosure & Precision or exponent & Cross-multiplied identities,
growth $Y_n\ge2^n-1$, and the named width bound.\\
\bottomrule
\end{longtable}

For an example of the Boolean translation, let
$F(t)=\prod_{i<t}\chi(i)$ with $\chi(i)\in\{0,1\}$.
The implication $F(t)=1\Rightarrow\chi(i)=1$ for $i<t$ is proved by
quantifier-free induction with $i$ as a free parameter. Conversely, a
bounded scan storing the first zero proves, again by a quantifier-free
invariant, that $F(t)=0$ supplies an index $i<t$ with $\chi(i)=0$.
If every bounded instance equals one, this alternative is contradictory;
the decidability of $F(t)$ gives $F(t)=1$. Thus translating a finite
universal assertion does not quietly grant induction on that assertion
in a stronger language.

For the descending algorithms, the loop itself is a primitive recursion
with a supplied fuel value. Completeness is proved from the residual
coordinate bound, so no induction on a statement of the form ``all smaller
solutions are classified'' is needed. For the Legendre argument, the
only additional case distinctions are signs of computed integers and
whether a determinant is zero.

These details establish a finite syntactic route from the displayed
mathematics to iPRA proofs. They do not assert that the correctness of
every primitive recursive program follows automatically from its
definition: the specific preservation identities and decreasing measures
proved in this article are the required proof content.

\subsection{The constructive content of the quantifiers}

Pell existence has a particularly strong Skolemized form. There are fixed
primitive recursive $A(D),B(D)$ such that iPRA proves the universal closure
of the quantifier-free implication
\begin{equation}\label{eq:skolem}
\begin{split}
 D\ge2\ \wedge\ \isqrt(D)^2\ne D\quad\Longrightarrow\quad
 &0<B(D)\le\CFB(D),\\[-2pt]
 &A(D)\le\CFB(D),\\[-2pt]
 &A(D)^2=DB(D)^2+1.
\end{split}
\end{equation}
The last line is written additively so it has a literal natural-number
interpretation. The universal quantifier is introduced after the
quantifier-free induction arguments; it is not part of an induction formula.
Existential introduction then gives the ordinary existence theorem.

For negative Pell, the Boolean oddness of the computed period produces
either the explicit negative solution or a proof excluding any proposed
integer pair. For generalized Pell, the Boolean nonemptiness of the computed
seed list similarly gives either a seed or the implication that every
proposed solution would descend to a nonexistent seed. Thus the negative
answers are also proved constructively. No double-negation extraction,
Markov principle, excluded middle for arbitrary formulas, choice principle,
infinite pigeonhole principle, or real completeness principle is needed.

\section{A clear path to a complete Lean formalization}\label{sec:lean}

\subsection{Repository revision and inspected foundations}

The ProveIt snapshot inspected for this work is
\begin{center}
\small\ttfamily f1b0582aada7262b10924ea1ba12e2e490d2934b.
\end{center}
This pins the discussion to a concrete revision rather than to the moving
\code{main} branch. The main foundation is
\path{Logic/PrimitiveRecursiveArithmetic/ConciseIPRA/Lean/ConciseIPRA/}.
Its actual source, not just its project description, provides the following
objects \cite{proveit-pr,proveit-native}:
\begin{itemize}
\item \code{PR.lean}: the syntax \code{PR} of primitive recursive function
symbols, the term language \code{PTm}, the theory \code{PRAAx}, and the
derived rule \code{PRA.qfInd}. The induction constructor requires a proof
that its formula is quantifier-free.
\item \code{TeqReflect.lean}: characteristic-term reflection, including
\code{chi}, \code{chi\_iff}, and the derived quantifier-free induction rule
\code{Teq.qfInd}.
\item \code{TeqRec.lean}: the pair-accumulator implementation of primitive
recursion, useful for the several registers of a bounded algorithm.
\item \code{NativeBiInterpretation.lean}: the exact faithful translation
\code{pra\_iff\_native} from embedded iPRA derivability to native scalar
derivability, and its converse.
\end{itemize}
The pinned repository toolchain and Mathlib tag are both version
\code{v4.32.0}. The accompanying roadmap records the inspected blob SHAs,
full paths, and API names for reproducibility.

We also inspected the draft directory listing. It contains other recent
arithmetic and analysis packages, but no theorem from those ZIP archives is
assumed here. The argument is self-contained and needs no extraction or
trust in an unpublished proof inside another archive.

\subsection{Useful existing Pell files and their limits}

The repository file
\path{Computability/HilbertTenthProblem/Lean/Diophantine/Common/Pell.lean}
develops the special discriminants $a^2-1$, importing Mathlib's
\code{PellMatiyasevic}. Its recurrence, growth, and congruence formulas are
useful comparison material. This does not by itself prove existence for
every nonsquare $D$ in the embedded arithmetic theory.

The adjacent \code{PellMod.lean} contains a theorem
\code{exists\_dvd\_yn} whose proof uses
\code{Finite.exists\_ne\_map\_eq\_of\_infinite}. For the present
formalization, \cref{thm:modular} supplies the finite replacement with
$h^2+1$ states and a computed return exponent.

Similarly,
\path{Computability/DiophantineConsistency/Lean/DiophantineConsistency/NativePellBootstrap.lean}
uses ordinary Lean/Mathlib Pell propositions. A filename containing
\code{Native} does not make its theorems derivations in \code{PRAAx}.
The theorem's type determines that distinction.

Mathlib's \code{NumberTheory.Pell} supplies an ordinary Lean existence
theorem and a fundamental-unit classification; its documentation lists
continued fractions and the negative equation as further development
topics at the inspected source \cite{mathlib-pell}. Those host-theory
results can check mathematical compatibility, but importing them is not
a proof that iPRA derives \eqref{eq:skolem}.

\subsection{Two layers, with an explicit bridge}

The first layer proves the arithmetic algorithms and identities as ordinary
Lean statements about \code{Nat}, \code{Int}, and finite records.
The included \code{PellCore.lean} provides a starting point: quadratic-pair
multiplication and norm, associativity and conjugation, power recursion,
norm preservation, power addition, and the numerical bounds and solver
specification.

The second layer constructs the actual \code{PR} syntax for the fixed
functions, constructs a formula expressing \eqref{eq:skolem}, and produces
a proof term in the repository's derivability relation. The target has the
following shape; the formula name here denotes work still to be completed:
\begin{lstlisting}
-- PellCorrect : Fm PTm 1 is a quantifier-free formula.
-- It mentions explicit PR function symbols for the solver and bound.
-- Required final target:
-- ([] : List (Fm PTm 0)) |- [PRAAx] (Fm.all PellCorrect)
\end{lstlisting}
The ASCII turnstile in this display is explanatory notation. The supplied
\code{Transfer.lean} uses the repository's actual scoped turnstile notation
and the exact theorem \code{pra\_iff\_native}. Once an embedded derivation
$h$ exists, its translation is
\begin{lstlisting}
  (pra_iff_native).mp h
\end{lstlisting}
with simplification of the empty context. This step does not create the
missing embedded arithmetic proof; it transports that proof after it has
been constructed.

\subsection{Implementation modules and precise proof obligations}

The following order is a concrete completion plan.
\begin{enumerate}
\item \textbf{Arithmetic coding.} Implement signed integers, rational pairs,
finite vectors, total division, square-root search, bounded selectors, and
fixed-fuel Euclid as \code{PR} terms. Prove their defining laws and numerical
invariants using \code{PRA.qfInd}. Record the translation between these
codes and the ordinary Lean structures.
\item \textbf{Quadratic order.} Implement the sign test from
\eqref{eq:signformula}. Prove the rational-root lemma and the positive-
coefficient fraction witnesses. Derive the ordered-field calculations used
for the error estimates. All sign splits are decidable arithmetic splits.
\item \textbf{Finite collision.} Formalize the occupancy scan with its
success flag and stored indices. Compile the occupancy-count invariant to
an equation of a Boolean characteristic term. This one module supports
state periods, residue periods, and the alternative existence proof.
\item \textbf{Reduced states.} Encode the triple $(m,d,c)$ and prove
\cref{lem:reduced}. The exact division and the predecessor reconstruction
are the central obligations. Use the state count $S$ and bounded collision
to produce $\ell$ and prove \eqref{eq:unitdenom}.
\item \textbf{Convergents.} Store the last two numerator and denominator
values. Prove \eqref{eq:inv1}--\eqref{eq:det} by one simultaneous
quantifier-free induction, then derive \eqref{eq:cfnorm} and
\eqref{eq:convbound} by polynomial identities.
\item \textbf{Minimality and negative parity.} Implement denominator
bracketing with fuel $q+1$. The integral inverse of the consecutive-
convergent matrix is explicit, so the best-approximation and Legendre
proofs reduce to sign cases and determinant estimates. Combine these with
the period characterization to prove \cref{thm:parity} and
\eqref{eq:skolem}.
\item \textbf{Structure.} Formalize the three square-difference lemmas
for each descent. The fixed-fuel reconstruction invariant proves the
classification, negative square-root theorem, and generalized finite seed
theorem. No well-founded recursion is required in the object language.
\item \textbf{Arithmetic consequences.} Formalize the bounded Euclidean
trace for strong divisibility and the $h^2$ residue return. Encode rational
enclosures by cross products and their precision by the function $k+1$.
\item \textbf{Native translation and checking.} Apply the existing
\code{pra\_iff\_native} theorem. Compile the finished project, inspect
\code{\#print axioms} for the final declarations, and retain fresh build
logs. If generating a standalone checker certificate is added, replay it
with the appropriate ConciseIPRA checker as a separate final step.
\end{enumerate}

The alternative norm-collision proof can be formalized after the finite
collision module and before continued fractions. It provides an independent
route to general existence with bound \eqref{eq:H}, and can be useful while
the faster continued-fraction modules are under construction.

At the host layer, Lean's \code{ring} tactic can discharge polynomial
identities. For embedded iPRA, one must instead construct the corresponding
arithmetic derivations from the defining equations and proved ring laws,
or implement and prove sound a suitable reflection procedure. A successful
host-level \code{ring} call is not itself the embedded derivation. The
few displayed identities can also be expanded by direct rewrite proofs,
which avoids making a general reflection engine a prerequisite.

\subsection{What is supplied and what is not claimed}

The included Lean sources contain no admitted theorem placeholders and no
added axioms, but they have not been compiled: no Lean or Lake executable
was available in the delivery environment. Accordingly, even those short
source proofs are supplied as an unverified starting point. The complete
embedded iPRA derivation, a machine-checked general Pell theorem, and a
native proof-checker certificate remain formalization work. The human
proofs and the explicit module plan are the requested non-formalized
argument and outline, respectively.

The ordinary Lean starting point can be checked in the supplied standalone
project with the pinned toolchain:
\begin{lstlisting}
lake update
lake exe cache get
lake env lean PellCore.lean
\end{lstlisting}
The transfer example is checked inside the pinned ProveIt workspace after
building \code{ConciseIPRA.NativeBiInterpretation}; it is intentionally not
a standalone Lake target. A separate host axiom audit and an object-theory
derivation audit answer different questions and should both be retained
when the formalization is completed.

\section{Executable validation and reproducibility}\label{sec:validation}

The archive includes a Python standard-library implementation with exact
integer arithmetic, a deterministic test suite, seven worked
continued-fraction traces, Lean source outlines, and a provenance record.
All source-level loops have a stated finite range. An early return or break
is an optimization of a loop that holds its state fixed after success.
The integer-square-root implementation uses a finite binary-digit scan;
its results are checked against the square-root inequalities.

The executed suite passed all 12 test methods, with zero failures and
zero errors. The saved \code{test-results.json} records the scope:
\begin{center}
\small
\begin{tabular}{lr}
\toprule
Check & Number checked\\
\midrule
Nonsquare discriminants $2\le D\le500$ & 478\\
Reduced states over their first periods & 3\,879\\
Convergent norm identities & 7\,758\\
Independent bounded minimality cases & 116\\
Denominators in those minimality scans & 19\,212\\
Finite Dirichlet bin certificates & 750\\
Equal-signature integral quotients & 5\\
Norm-one descent reconstructions & 72\\
Negative descent reconstructions & 30\\
Generalized norm families & 175\\
Generalized seed-power reductions & 772\\
Independently enumerated generalized solutions & 232\\
Finite modular periods & 180\\
Strong divisibility comparisons & 1\,014\\
Rational enclosures & 77\\
\bottomrule
\end{tabular}
\end{center}

These finite computations support the implementation and check indexing
and sign conventions. They are not substitutes for the universal proofs,
and the JSON traces are arithmetic certificates rather than derivations in
the embedded iPRA checker. The enormous two-collision existence algorithm
is not run in full; its finite Dirichlet component and quotient identities
are exercised separately, while the efficient continued-fraction algorithm
is run over the stated discriminant range.

To reproduce the executable checks from the archive's \code{code} directory:
\begin{lstlisting}
python3 test_pell.py
python3 pell_reference.py
python3 export_certificates.py
\end{lstlisting}
To rebuild this PDF, run \code{python3 build.py} from the archive root;
it invokes \code{pdflatex} repeatedly to resolve references. The manifest
records SHA-256 hashes of the delivered files. The source material reused
from ProveIt is identified by repository revision and file blob SHA in
\code{lean/ROADMAP.md} and \code{evidence/provenance.json}.

\section{Conclusion}

Pell's equation admits a particularly direct development in the requested
weak constructive arithmetic. The apparent infinitary steps have finite
arithmetic replacements: a finite state space forces the continued-fraction
period; a unimodular two-vector calculation gives Legendre's criterion;
and integer inverse multiplication gives the complete solution
classification with a fixed descent budget. The same unit reduces every
generalized norm equation to a finite seed scan.

The resulting proofs supply explicit primitive recursive witnesses for
existence, periods, exponents, seed reductions, and approximation precision.
They also give constructive negative answers when negative or generalized
Pell has no solution. The formalization route preserves this content by
targeting an actual iPRA derivation and only then applying ProveIt's native
translation.

\begin{thebibliography}{9}
\raggedright
\bibitem{conrad1}
Keith Conrad, \emph{Pell's Equation, I}, author-hosted expository notes.
\url{https://kconrad.math.uconn.edu/blurbs/ugradnumthy/pelleqn1.pdf}.

\bibitem{conrad2}
Keith Conrad, \emph{Pell's Equation, II}, author-hosted expository notes.
\url{https://kconrad.math.uconn.edu/blurbs/ugradnumthy/pelleqn2.pdf}.

\bibitem{lynn}
Ben Lynn, \emph{Continued Fractions: The Pell Equation}, Stanford-hosted notes.
\url{https://crypto.stanford.edu/pbc/notes/contfrac/pell.html}.

\bibitem{mathlib-pell}
The Mathlib community, \emph{Mathlib.NumberTheory.Pell}, official
documentation and linked source.
\url{https://leanprover-community.github.io/mathlib4_docs/Mathlib/NumberTheory/Pell.html}.

\bibitem{proveit-pr}
Vladimir Reshetnikov et al., \emph{ProveIt: ConciseIPRA, PR.lean},
revision \texttt{f1b0582aada7262b10924ea1ba12e2e490d2934b}.
\url{https://github.com/VladimirReshetnikov/ProveIt/blob/f1b0582aada7262b10924ea1ba12e2e490d2934b/Logic/PrimitiveRecursiveArithmetic/ConciseIPRA/Lean/ConciseIPRA/PR.lean}.

\bibitem{proveit-native}
Vladimir Reshetnikov et al., \emph{ProveIt: ConciseIPRA,
NativeBiInterpretation.lean}, same revision.
\url{https://github.com/VladimirReshetnikov/ProveIt/blob/f1b0582aada7262b10924ea1ba12e2e490d2934b/Logic/PrimitiveRecursiveArithmetic/ConciseIPRA/Lean/ConciseIPRA/NativeBiInterpretation.lean}.

\bibitem{pellpage}
\emph{Pell's equation}, the encyclopedia page specified in the task;
used for orientation and scope.
\url{https://en.wikipedia.org/wiki/Pell\%27s_equation}.
\end{thebibliography}

\end{document}



